%%%%%%%%%%  scan idx 284  |  folio 271  %%%%%%%%%%
\SGAPrintPage{271}%
\(y\mapsto\operatorname{rg}_{k(y)}H^p(E\otimes_X^{\mathrm L}k(y))\) is locally constant by (5.5 (ii)). The constructibility assertion of (i) follows at once. The other assertions follow from (5.5 (ii) and (iii)) and (5.2 (iv)), except for the converse of (ii), for which it suffices, after reducing to \(E\) strictly perfect (by localization and truncation), to apply (EGA III 7.6.2 and 7.6.5) to the functor \(T_{-p}=H^p(E\otimes\,.\,)\).

\sgaStableTarget{sga6.e03.prop.5-7}\textbf{Proposition 5.7.} (EGA III 7.7.5 \uline{and} 7.9.4). \uline{Let} \(f:X\to Y\) \uline{be a proper morphism of quasi-compact schemes}, \(E\in\Ob D^b(f)_{\uparf}\) (4.2), \uline{and} \(p\in\mathbb Z\). \uline{Then}, \uline{subject to the validity of} (2.1), \uline{hence for example if} \(f\) \uline{is projective or if} \(Y\) \uline{is noetherian}:
\begin{enumerate}[label=(\roman*)]
\item \uline{For} \(y\in Y\), \uline{the} \(H^i(Rf_*(E)\otimes_Y^{\mathrm L}k(y))\) \uline{are vector spaces of finite rank over} \(k(y)\), \uline{all but finitely many zero}, \uline{and the function}
\[
 y\longmapsto\sum(-1)^i\operatorname{rg}_{k(y)}H^i(Rf_*(E)\otimes_Y^{\mathrm L}k(y))
\]
\uline{is locally constant on} \(Y\).
\item \uline{The function}
\[
 y\longmapsto\operatorname{rg}_{k(y)}H^p(Rf_*(E)\otimes_Y^{\mathrm L}k(y))
\]
\uline{is constructible and upper semi-continuous on} \(Y\).
\item \uline{Suppose} \(Y\) \uline{quasi-separated}. \uline{Then}, \uline{if the functor} \(M\mapsto R^pf_*(E\otimes_Y^{\mathrm L}M)\) \uline{is exact on the category of quasi-coherent} \(\underline{O}_Y\)-\uline{Modules}, \(R^pf_*(E)\) \uline{is a locally free} \(\underline{O}_Y\)-\uline{Module of finite type}, \uline{and the function of} (ii) \uline{is locally constant}. \uline{Conversely}, \uline{if} \(Y\) \uline{is noetherian and reduced}, \uline{and if the function of} (ii) \uline{is locally constant}, \uline{the functor} \(M\mapsto R^pf_*(E\otimes_Y^{\mathrm L}M)\) \uline{is exact on the category of quasi-coherent} \(\underline{O}_Y\)-\uline{Modules}.
\end{enumerate}

%%%%%%%%%%  scan idx 285  |  folio 272  %%%%%%%%%%
\SGAPrintPage{272}%
\uline{Proof}. By (4.8), \(Rf_*(E)\) is a perfect complex.

Assertion (i) (resp. (ii)) therefore follows from (5.5 (i)) (resp. (5.6 (i)) ). Finally, if \(M\) is a quasi-coherent \(\underline{O}_Y\)-Module, the canonical arrow
\[
 Rf_*(E)\otimes_Y^{\mathrm L}M\longrightarrow Rf_*(E\otimes_Y^{\mathrm L}M)
\]
is an isomorphism by (3.7), so assertion (iii) follows from (5.6 (ii)).

\sgaStableTarget{sga6.e03.rem.5-8}\textbf{Remarks 5.8.}
\begin{enumerate}[label=\alph*)]
\item In the situation of (5.7), denote by \(i_y:X_y\to X\) the inclusion of the fiber of \(X\) over a point \(y\) of \(Y\). One can show (cf. (EGA III 6) and (IV 3.1.0)) that, for \(f\) flat (resp. \(E\) bounded with quasi-coherent and \(Y\)-flat components), we have the Künneth formula:
\[
 H^i(Rf_*(E)\otimes_Y^{\mathrm L}k(y))\xrightarrow{\sim}H^i(X_y,E_y),
\]
where \(E_y=Li_y^*(E)\) (resp. \(i_y^*(E)\)).
%%%%%%%%%%  scan idx 286  |  folio 273  %%%%%%%%%%
\SGAPrintPage{273}%
\item One can develop for complex analytic spaces a body of results analogous to that of this section (exchange, semi-continuity, continuity properties)\SGAAuthorFootnote{1}{It was moreover in this setting that these properties were first formulated (see for example Grauert [1]), before being transposed to Algebraic Geometry in (EGA III 7).}. The essential ingredients are: 1) the finiteness theorem for proper morphisms (``Grauert's theorem''); 2) an analogue of the projection formula (3.7), which may be stated as follows: let \(f:X\to Y\) be a morphism of analytic spaces, where \(X\) is supposed of bounded dimension, and let \(E\in\Ob D^-(X)\), \(M\in\Ob D^-(Y)_{\coh}\) (``coh'' means ``with coherent cohomology''). The canonical arrow
\[
 Rf_*(E)\otimes^{\mathrm L}M\longrightarrow Rf_*(E\otimes_Y^{\mathrm L}M)
\]
is an isomorphism (the hypothesis on \(X\) ensures (T.F II 5.12) that \(f_*\) has finite cohomological dimension, hence that \(Rf_*\) is defined on \(D(X)\)).

Unlike statement 1), statement 2) is essentially trivial, and follows from a dévissage analogous to that of the proof of (3.7); we moreover have greater freedom since \(Rf_*\) is defined on \(D(X)\), and in particular no hypothesis on \(E\) is needed (apart from \(E\in\Ob D^-(X)\)).
\end{enumerate}

\sgaStableTarget{sga6.e03.bibliography}\subsection*{BIBLIOGRAPHY}
\begin{enumerate}[label={[\arabic*]}]
\item[{[1]}] Grauert, H., Ein Theorem der Analytischen Garbentheorie \ldots, Pub. IHES, n\textsuperscript{o}5.
\item[{[H]}] Hartshorne, R., Residues and Duality, Lecture Notes n\textsuperscript{o}20, Springer, 1966.
\item[{[T.F]}] Godement, R., Théorie des faisceaux, Hermann, 1958.
\end{enumerate}

% SGA_NATIVE_CONTINUATION_273
\providecommand{\Ob}{\operatorname{Ob}}
\providecommand{\Mod}{\operatorname{Mod}}
\providecommand{\Coh}{\operatorname{Coh}}
\providecommand{\Loclib}{\operatorname{Loclib}}
\providecommand{\coh}{\mathrm{coh}}
\providecommand{\parf}{\mathrm{parf}}
\providecommand{\naif}{\mathrm{naif}}
\providecommand{\loc}{\mathrm{loc}}
\providecommand{\fin}{\mathrm{fin}}
\providecommand{\clA}{\operatorname{cl}}
\providecommand{\Rg}{\operatorname{rg}}
\providecommand{\varinjlim}{\varinjlim}
\providecommand{\varprojlim}{\varprojlim}
\providecommand{\id}{\operatorname{id}}
\providecommand{\Th}{\operatorname{Th}}
\providecommand{\Hom}{\operatorname{Hom}}
\providecommand{\Dcat}{\mathrm D}
\providecommand{\Kcat}{\mathrm K}
\providecommand{\Kdot}{K_{\bullet}}
\providecommand{\Kprime}{K'}
\providecommand{\naif}{\mathrm{na\ddot{\i}f}}


% SGA_REVIEWED_FRENCH_CONTINUATION_273
% Exposé IV macros
\renewcommand{\Ob}{\operatorname{ob}}
\providecommand{\Mod}{\operatorname{Mod}}
\providecommand{\Coh}{\operatorname{Coh}}
\providecommand{\Loclib}{\operatorname{Loclib}}
\providecommand{\coh}{\mathrm{coh}}
\providecommand{\parf}{\mathrm{parf}}
\providecommand{\naif}{\text{naïf}}
\providecommand{\loc}{\mathrm{loc}}
\providecommand{\fin}{\mathrm{fin}}
\providecommand{\clA}{\operatorname{cl}}
\providecommand{\Rg}{\operatorname{rg}}
\providecommand{\varinjlim}{\varinjlim}
\providecommand{\varprojlim}{\varprojlim}
\renewcommand{\id}{\operatorname{id}}
\providecommand{\Th}{\operatorname{Th}}
\providecommand{\Hom}{\operatorname{Hom}}
\providecommand{\Dcat}{\mathrm D}
\providecommand{\Kcat}{\mathrm K}
\providecommand{\Kdot}{K_{\bullet}}
\providecommand{\Kprime}{K'}
\providecommand{\naif}{\text{naïf}}



% SGA_SOURCE_MACRO_RESTORATION_273: native provide must not shadow the reviewed accented label.
\renewcommand{\naif}{\text{naïf}}

